
\documentclass{article}
\usepackage{colortbl}
\usepackage{makecell}
\usepackage{multirow}
\usepackage{supertabular}

\begin{document}

\newcounter{utterance}

\centering \large Interaction Transcript for game `codenames', experiment `concrete', episode 9 with reflexion{-}Qwen3.5{-}9B.
\vspace{24pt}

{ \footnotesize  \setcounter{utterance}{1}
\setlength{\tabcolsep}{0pt}
\begin{supertabular}{c@{$\;$}|p{.15\linewidth}@{}p{.15\linewidth}p{.15\linewidth}p{.15\linewidth}p{.15\linewidth}p{.15\linewidth}}
    \# & $\;$A & \multicolumn{4}{c}{Game Master} & $\;\:$B\\
    \hline

    \theutterance \stepcounter{utterance}  
    & & & \multicolumn{2}{p{0.3\linewidth}}{
        \cellcolor[rgb]{0.9,0.9,0.9}{
            \makecell[{{p{\linewidth}}}]{
                \texttt{\tiny{[GM$|$GM]}}
                {'hidden': {'team': ['salad', 'panther'], 'innocent': ['cranberry', 'dustpan', 'longbow', 'toothbrush', 'warthog'], 'opponent': [], 'assassin': ['eagle']}, 'revealed': {'team': {'team': [], 'innocent': [], 'opponent': [], 'assassin': []}, 'opponent': {'team': [], 'innocent': [], 'opponent': [], 'assassin': []}}}
            }
        }
    }
    & & \\ \\

    \theutterance \stepcounter{utterance}  
    & & \multicolumn{4}{p{0.6\linewidth}}{
        \cellcolor[rgb]{0.9,0.9,0.9}{
            \makecell[{{p{\linewidth}}}]{
                \texttt{\tiny{[P1$\langle$GM]}}
                \texttt{You have attempted this same task before and did not succeed. These are the plans you wrote after those attempts.} \\
\texttt{Trial \#1: New plan: I failed because I included "salad" in the "clothing" clue, which caused my partner to guess "jockstrap" (correct) but then likely guessed "salad" as a distractor or misinterpretation, ending the turn prematurely before I could clear the final word. More critically, in the final turn with only "salad" and "panther" remaining, I chose "leafy" for "salad", but "panther" was left unguessed. I should have prioritized grouping the last two words under a single, strong clue like "big cat" (for panther) or "green vegetable" (for salad) to ensure both are guessed in one turn, or simply guessed the remaining words directly if the game allowed, but since I must provide a clue, I need to ensure the clue is specific enough to trigger the guess of the *only* remaining target without ambiguity. Actually, looking back, the loss likely happened because "salad" was guessed as a distractor or the turn ended without clearing "panther". The core mistake was not ensuring the final clue covers *all* remaining targets effectively. Next time, with only two words left like "salad" and "panther", I will use a clue that strongly evokes one while avoiding the other, or better yet, realize that if I can't group them, I must pick the one that is most distinct. Wait, the rule says "guess all... before the other team". If I clue "salad" and they guess "salad", "panther" remains. If they guess "panther" on "leafy" (wrong), they might guess a distractor. The fatal error was leaving "panther" unguessed in the final turn. I should have used a clue like "feline" for "panther" alone, or if I could group, "food and animal" (but clue must be one word). Since I can't group unrelated words easily, I should have saved "panther" for a dedicated turn or chosen a clue that forces the guess of the harder word first. Actually, the simplest fix: In the final turn with "salad" and "panther", I should have guessed "panther" with a clue like "big cat" because "salad" is so obvious it might be guessed by mistake or the turn ends. No, the goal is to guess *all*. If I clue "big cat" {-}> "panther" guessed. Then "salad" remains. I need to win in one turn if possible. But I can't clue two unrelated things. The mistake was not realizing I couldn't clear both in one go and thus failed to plan for the second turn. Wait, the game ended with LOSE. This implies the opponent cleared their words or I failed to clear mine. With "salad" and "panther} \\
\\ 
\texttt{Trial \#2: New plan: I failed because I used "feline" for "panther," which is a strong clue, but I left "salad" unclued, allowing the opponent to clear their words or the game to end before I could guess "salad." The core mistake was not realizing that unrelated words like "salad" and "panther" cannot be grouped under a single semantic clue without ambiguity or risk. Since I must clear all words before the opponent, I should have prioritized the harder word ("panther") in the first turn with a specific clue like "big cat," and then immediately followed up with a dedicated turn for "salad" using a clue like "green vegetable" or "lettuce mix." In the final turn, I must ensure the clue targets the *only* remaining word to avoid leaving any targets unguessed. Next time, I will always verify that my clue covers all remaining targets or plan a strict sequence of single{-}word clues to guarantee no word is left behind.} \\
\\ 
\\ 
\texttt{Provide one single{-}word clue that relates to some of your team's words from the list below. You can choose to how many and to which words your clue relates to.} \\
\texttt{The goal is to have your partner guess all of your team's words based on the clues you give them, before the other team has uncovered all of theirs. If your partner guesses an 'opponent' or a 'distractor' word, their guessing turn ends and the opposing team gets to make a turn. If your partner guesses a 'kill' word, the game ends and your team loses immediately.} \\
\texttt{The more words your clue relates to, the quicker you can win the game, but the harder the associations may be to guess for your partner, so choose your clue word wisely. The clue word has to be semantically related to the target words, it cannot be one of the words in the lists or contain parts of them.} \\
\\ 
\texttt{Always give your single{-}word clue and your comma{-}separated list of related target words in the following format and make your answers as short as possible, never include any other text than is required in this form:} \\
\\ 
\texttt{CLUE: <WORD>} \\
\texttt{TARGETS: <TARGETS>} \\
\\ 
\texttt{Your team words are: salad, panther.} \\
\texttt{Your opponent's team words are: .} \\
\texttt{Distractor words are: cranberry, dustpan, longbow, toothbrush, warthog.} \\
\texttt{The kill words are: eagle.} \\
            }
        }
    }
    & \\ \\

    \theutterance \stepcounter{utterance}  
    & \multicolumn{4}{p{0.6\linewidth}}{
        \cellcolor[rgb]{0.9,0.9,0.9}{
            \makecell[{{p{\linewidth}}}]{
                \texttt{\tiny{[P1$\rangle$GM]}}
                \texttt{CLUE: feline} \\
\texttt{TARGETS: panther} \\
            }
        }
    }
    & & \\ \\

    \theutterance \stepcounter{utterance}  
    & & & \multicolumn{2}{p{0.3\linewidth}}{
        \cellcolor[rgb]{0.9,0.9,0.9}{
            \makecell[{{p{\linewidth}}}]{
                \texttt{\tiny{[GM$|$GM]}}
                \texttt{feline} \\
            }
        }
    }
    & & \\ \\

    \theutterance \stepcounter{utterance}  
    & & & \multicolumn{2}{p{0.3\linewidth}}{
        \cellcolor[rgb]{0.9,0.9,0.9}{
            \makecell[{{p{\linewidth}}}]{
                \texttt{\tiny{[GM$|$GM]}}
                ['panther']
            }
        }
    }
    & & \\ \\

    \theutterance \stepcounter{utterance}  
    & & & \multicolumn{2}{p{0.3\linewidth}}{
        \cellcolor[rgb]{0.9,0.9,0.9}{
            \makecell[{{p{\linewidth}}}]{
                \texttt{\tiny{[GM$|$GM]}}
                {'word': 'panther', 'assignment': 'team'}
            }
        }
    }
    & & \\ \\

    \theutterance \stepcounter{utterance}  
    & & & \multicolumn{2}{p{0.3\linewidth}}{
        \cellcolor[rgb]{0.9,0.9,0.9}{
            \makecell[{{p{\linewidth}}}]{
                \texttt{\tiny{[GM$|$GM]}}
                \texttt{opponent has won} \\
            }
        }
    }
    & & \\ \\

    \theutterance \stepcounter{utterance}  
    & \multicolumn{4}{p{0.6\linewidth}}{
        \cellcolor[rgb]{0.95,0.95,0.95}{
            \makecell[{{p{\linewidth}}}]{
                \texttt{\tiny{[P2$\rightarrow$P1]}}
                \texttt{Post{-}episode Reflexion (not sent to GM):} \\
\texttt{New plan: I failed because I only clued "panther" with "feline," leaving "salad" unclued, which allowed the opponent to clear their words first. The core mistake was not realizing that unrelated words like "salad" and "panther" cannot be grouped under a single semantic clue without risking ambiguity or leaving one behind. Since I must clear all words before the opponent, I should have prioritized the harder word ("panther") in the first turn with a specific clue like "big cat," and then immediately followed up with a dedicated turn for "salad" using a clue like "green vegetable" or "lettuce mix." In the final turn, I must ensure the clue targets the *only* remaining word to avoid leaving any targets unguessed. Next time, I will always verify that my clue covers all remaining targets or plan a strict sequence of single{-}word clues to guarantee no word is left behind.} \\
            }
        }
    }
    & & \\ \\

\end{supertabular}
}

\end{document}
